\addcontentsline{toc}{subsubsection}{Solution 6. The Sum--Product Squeeze}

\begin{dlsummary}[title={\bfseries Solution 6. The Sum--Product Squeeze},breakable]

\begin{proof}
Set
\[
g(x)=\frac1{f(x)},\qquad k=\frac{g(2)}2.
\]
The given inequalities are equivalent to
\begin{align}
g(x+y)&\ge k(g(x)+g(y)), \tag{1}\\
g(xy)&\le k\,g(x)g(y). \tag{2}
\end{align}

\textbf{Step 1.} Setting $x=y=1$ gives
\[
g(1)\le1,\qquad kg(1)\ge1,
\]
so $k\ge1$.

If $k>1$, repeated application of (1), with
increments $1/n$, yields
\[
g(2)\ge k^n g(1).
\]
Letting $n\to\infty$ gives a contradiction.
Thus $k=1$, and consequently
\[
g(1)=1,\qquad g(2)=2.
\]

\textbf{Step 2.} By (1) and (2),
\[
2g(x)\le g(2x)\le g(2)g(x)=2g(x).
\]
Hence
\[
\boxed{g(2x)=2g(x).}
\]
In particular, $g(2^n)=2^n$ for every
positive integer $n$.

\textbf{Step 3.} As $g$ is positive and
superadditive, it is strictly increasing.

Fix $x>0$ and put
\[
m_n=\left\lfloor\frac{2^n}{x}\right\rfloor.
\]
For all sufficiently large $n$,
\[
2^n=g(2^n)\ge g(m_nx)\ge m_ng(x).
\]
Therefore,
\[
g(x)\le
\frac{2^n}{\lfloor 2^n/x\rfloor}
\xrightarrow[n\to\infty]{}x.
\]
Thus $g(x)\le x$ for every $x>0$.

Conversely, by (2),
\[
1=g(1)\le g(x)g(1/x)\le\frac{g(x)}x,
\]
so $g(x)\ge x$.

We conclude that $g(x)=x$, whence
\[
\boxed{f(x)=\frac1x}.
\]

Finally, direct substitution gives
\[
f(x+y)(f(x)+f(y))
=2f(2)f(x)f(y)
=f(xy)=\frac1{xy}.
\]
Thus the proposed function satisfies
both inequalities.
\end{proof}

\end{dlsummary}